% GENERATED FILE. Edit canonical lessons, then run npm run generate:books.
% canonical-source-hash: fnv1a64:7eee2cef1d600d7a
% canonical-lessons: BN-C360-boktrita, BN-C360-bokta, BN-C360-sobhapoti, BN-C360-uposthito, BN-C360-sombodhon-kora

\chapter{Speeches and Talks}
\label{ch:bn-speeches-and-talks}
\hlchaptermodality{360}

\begin{chapteropening}
\textbf{By the end of this chapter:} \emph{I can talk about a speech, the speaker and who was there.}

Complete the last lesson of chapter 360: I can talk about a speech, the speaker and who was there.
\end{chapteropening}

\section[\texorpdfstring{\bn{বক্তৃতা} (bôktritā)}{বক্তৃতা (bôktritā)}]{\bn{বক্তৃতা} (bôktritā) — a speech; a lecture}

\label{lesson:BN-C360-boktrita}



\begin{quote}
Before the new one: say the Bengali for to sidestep, then the Bengali for acceptable.
\end{quote}

\subsection*{You'll want to know: \bn{বক্তৃতা}}
\textbf{\bn{বক্তৃতা}} — \emph{bôktritā} — \textquotedblleft{}a speech; a lecture\textquotedblright{}.

\begin{quote}
\bn{কাল} \bn{কলেজে} \bn{একটা} \bn{বক্তৃতা} \bn{আছে।}
\end{quote}

\begin{quote}
\emph{kāl kôleje ekṭā bôktritā āchhe.}
\end{quote}

\begin{quote}
There is a lecture at the college tomorrow.
\end{quote}

\begin{cousinweb}
\textbf{\bn{বক্তৃতা}} (\emph{bôktritā}) is Sanskrit \emph{vaktṛtā}, the gift of speech, eloquence (Monier-Williams). Bengali uses it for the speech itself: a lecture or an address.
\end{cousinweb}

\begin{grammarlens}[title={What you've built}]
One more word for giving reasons and opinions.
\end{grammarlens}

\subsection*{Your turn}
\begin{itemize}
\raggedright
  \item \emph{Say it:} \emph{bôktritā}
  \item \emph{Say it:} \emph{bôktritā}, once more
  \item \emph{Say it:} \emph{kāl kôleje ekṭā bôktritā āchhe.} — There is a lecture at the college tomorrow.
  \item \emph{Recall:} say the Bengali for to sidestep, then the Bengali for acceptable, then say \emph{bôktritā} again
\end{itemize}

\begin{tcolorbox}[breakable,colback=teal!4,colframe=teal!35!black,title={Before you move on}]
What does \bn{বক্তৃতা} mean? (\textquotedblleft{}A speech; a lecture\textquotedblright{}.) Say it once more.
\end{tcolorbox}
\section[\texorpdfstring{\bn{বক্তা} (bôktā)}{বক্তা (bôktā)}]{\bn{বক্তা} (bôktā) — a speaker}

\label{lesson:BN-C360-bokta}



\begin{quote}
Before the new one: say the Bengali for acceptable, then the Bengali for a lecture.
\end{quote}

\subsection*{You'll want to know: \bn{বক্তা}}
\textbf{\bn{বক্তা}} — \emph{bôktā} — \textquotedblleft{}a speaker\textquotedblright{}.

\begin{quote}
\bn{সভার} \bn{প্রধান} \bn{বক্তা} \bn{একজন} \bn{বিজ্ঞানী।}
\end{quote}

\begin{quote}
\emph{sôbhār prôdhān bôktā ekjôn bigyānī.}
\end{quote}

\begin{quote}
The main speaker at the meeting is a scientist.
\end{quote}

\begin{cousinweb}
\textbf{\bn{বক্তা}} (\emph{bôktā}) is Sanskrit \emph{vaktṛ}, a speaker (Monier-Williams), in its nominative form vaktā.
\end{cousinweb}

\begin{grammarlens}[title={What you've built}]
One more word for giving reasons and opinions.
\end{grammarlens}

\subsection*{Your turn}
\begin{itemize}
\raggedright
  \item \emph{Say it:} \emph{bôktā}
  \item \emph{Say it:} \emph{bôktā}, once more
  \item \emph{Say it:} \emph{sôbhār prôdhān bôktā ekjôn bigyānī.} — The main speaker at the meeting is a scientist.
  \item \emph{Recall:} say the Bengali for acceptable, then the Bengali for a lecture, then say \emph{bôktā} again
\end{itemize}

\begin{tcolorbox}[breakable,colback=teal!4,colframe=teal!35!black,title={Before you move on}]
What does \bn{বক্তা} mean? (\textquotedblleft{}A speaker\textquotedblright{}.) Say it once more.
\end{tcolorbox}
\section[\texorpdfstring{\bn{সভাপতি} (sôbhāpôti)}{সভাপতি (sôbhāpôti)}]{\bn{সভাপতি} (sôbhāpôti) — a chairperson; the chair (of a meeting)}

\label{lesson:BN-C360-sobhapoti}



\begin{quote}
Before the new one: say the Bengali for a lecture, then the Bengali for a speaker.
\end{quote}

\subsection*{You'll want to know: \bn{সভাপতি}}
\textbf{\bn{সভাপতি}} — \emph{sôbhāpôti} — \textquotedblleft{}a chairperson; the chair (of a meeting)\textquotedblright{}.

\begin{quote}
\bn{সভাপতি} \bn{সভা} \bn{শুরু} \bn{করলেন।}
\end{quote}

\begin{quote}
\emph{sôbhāpôti sôbhā shuru kôrlen.}
\end{quote}

\begin{quote}
The chair opened the meeting.
\end{quote}

\begin{cousinweb}
\textbf{\bn{সভাপতি}} (\emph{sôbhāpôti}) is Sanskrit \emph{sabhāpati}, the head of an assembly (Monier-Williams), from \bn{সভা} with \emph{pati} 'master'.
\end{cousinweb}

\begin{grammarlens}[title={What you've built}]
One more word for giving reasons and opinions.
\end{grammarlens}

\subsection*{Your turn}
\begin{itemize}
\raggedright
  \item \emph{Say it:} \emph{sôbhāpôti}
  \item \emph{Say it:} \emph{sôbhāpôti}, once more
  \item \emph{Say it:} \emph{sôbhāpôti sôbhā shuru kôrlen.} — The chair opened the meeting.
  \item \emph{Recall:} say the Bengali for a lecture, then the Bengali for a speaker, then say \emph{sôbhāpôti} again
\end{itemize}

\begin{tcolorbox}[breakable,colback=teal!4,colframe=teal!35!black,title={Before you move on}]
What does \bn{সভাপতি} mean? (\textquotedblleft{}A chairperson; the chair (of a meeting)\textquotedblright{}.) Say it once more.
\end{tcolorbox}
\section[\texorpdfstring{\bn{উপস্থিত} (upôsthitô)}{উপস্থিত (upôsthitô)}]{\bn{উপস্থিত} (upôsthitô) — present; in attendance}

\label{lesson:BN-C360-uposthito}



\begin{quote}
Before the new one: say the Bengali for a speaker, then the Bengali for a chairperson.
\end{quote}

\subsection*{You'll want to know: \bn{উপস্থিত}}
\textbf{\bn{উপস্থিত}} — \emph{upôsthitô} — \textquotedblleft{}present; in attendance\textquotedblright{}.

\begin{quote}
\bn{সভায়} \bn{একশো} \bn{লোক} \bn{উপস্থিত} \bn{ছিল।}
\end{quote}

\begin{quote}
\emph{sôbhāy eksho lok upôsthitô chhilô.}
\end{quote}

\begin{quote}
A hundred people were at the meeting.
\end{quote}

\begin{cousinweb}
\textbf{\bn{উপস্থিত}} (\emph{upôsthitô}) is Sanskrit \emph{upasthita}, standing near, present (Monier-Williams). It is the word for being present at a meeting or in class.
\end{cousinweb}

\begin{grammarlens}[title={What you've built}]
One more word for giving reasons and opinions.
\end{grammarlens}

\subsection*{Your turn}
\begin{itemize}
\raggedright
  \item \emph{Say it:} \emph{upôsthitô}
  \item \emph{Say it:} \emph{upôsthitô}, once more
  \item \emph{Say it:} \emph{sôbhāy eksho lok upôsthitô chhilô.} — A hundred people were at the meeting.
  \item \emph{Recall:} say the Bengali for a speaker, then the Bengali for a chairperson, then say \emph{upôsthitô} again
\end{itemize}

\begin{tcolorbox}[breakable,colback=teal!4,colframe=teal!35!black,title={Before you move on}]
What does \bn{উপস্থিত} mean? (\textquotedblleft{}Present; in attendance\textquotedblright{}.) Say it once more.
\end{tcolorbox}
\section[\texorpdfstring{\bn{সম্বোধন} \bn{করা} (sômbodhôn kôrā)}{সম্বোধন করা (sômbodhôn kôrā)}]{\bn{সম্বোধন} \bn{করা} (sômbodhôn kôrā) — to address a crowd; to address (someone)}

\label{lesson:BN-C360-sombodhon-kora}



\begin{quote}
Before the new one: say the Bengali for a chairperson, then the Bengali for in attendance.
\end{quote}

\subsection*{You'll want to know: \bn{সম্বোধন} \bn{করা}}
\textbf{\bn{সম্বোধন} \bn{করা}} — \emph{sômbodhôn kôrā} — \textquotedblleft{}to address a crowd; to address (someone)\textquotedblright{}.

\begin{quote}
\bn{নেতা} \bn{মিছিলের} \bn{লোকদের} \bn{সম্বোধন} \bn{করলেন।}
\end{quote}

\begin{quote}
\emph{netā michhiler lokder sômbodhôn kôrlen.}
\end{quote}

\begin{quote}
The leader addressed the people in the procession.
\end{quote}

Past (he/she …): \bn{তিনি} \bn{সভাকে} \bn{সম্বোধন} \bn{করলেন}, \emph{tini sôbhāke sômbodhôn kôrlen}.

\begin{cousinweb}
\textbf{\bn{সম্বোধন} \bn{করা}} (\emph{sômbodhôn kôrā}) is Sanskrit \emph{sambodhana}, calling to someone (Monier-Williams), with \bn{করা}: to speak to a person or a crowd directly.
\end{cousinweb}

\begin{grammarlens}[title={What you've built}]
One more word for giving reasons and opinions.
\end{grammarlens}

\subsection*{Your turn}
\begin{itemize}
\raggedright
  \item \emph{Say it:} \emph{sômbodhôn kôrā}
  \item \emph{Say it:} \emph{sômbodhôn kôrā}, once more
  \item \emph{Say it:} \emph{netā michhiler lokder sômbodhôn kôrlen.} — The leader addressed the people in the procession.
  \item \emph{Recall:} say the Bengali for a chairperson, then the Bengali for in attendance, then say \emph{sômbodhôn kôrā} again
\end{itemize}

\begin{tcolorbox}[breakable,colback=teal!4,colframe=teal!35!black,title={Before you move on}]
What does \bn{সম্বোধন} \bn{করা} mean? (\textquotedblleft{}To address a crowd; to address (someone)\textquotedblright{}.) Say it once more.
\end{tcolorbox}
